\section{An evaluation plan}

The programme below is ordered by cost of evidence, not by ambition. Each
phase produces a result that stands alone, so the work degrades gracefully if
time runs out. Phase 0 is not optional: four of the six candidates depend on
the same annotated substrate, and building it once is the difference between
one paper and none.

\subsection{Phase 0: build the substrate}

Everything rests on one artifact, a role-annotated meeting slice. Use the AMI
scenario series, because it already assigns four fixed roles, project manager,
industrial designer, user interface designer and marketing expert, and because
its meetings come in linked series of four, which is what makes cross-meeting
contradiction items possible at all. Borrowing the roles from the corpus rather
than inventing them is what keeps the annotation defensible.

Replay the audio through the real capture path so that recognition and
diarization error are present at their measured rates. Then annotate, per 30
second window: a per-role relevance label, per-role reference extractions for
decisions, actions and questions, and a decision-closure timestamp marking the
moment after which surfacing an item is worthless. Add the abstention classes
explicitly, windows whose correct behaviour is discard, log silently, or
surface, plus context-gap queries whose true answer is that no source holds it.

Release this slice. It is plausibly a larger contribution than any single
experiment, because theme 2 showed that no existing benchmark can express
per-participant divergence at all.

\subsection{Phases 1 to 5, in order}

{\footnotesize
\begin{longtable}{@{}p{0.55cm}p{4.5cm}p{4.6cm}p{4.9cm}@{}}
\toprule
\textbf{Ph} & \textbf{Experiment} & \textbf{Headline metric} & \textbf{Why it is placed here} \\
\midrule
\endfirsthead
\toprule
\textbf{Ph} & \textbf{Experiment} & \textbf{Headline metric} & \textbf{Why it is placed here} \\
\midrule
\endhead
\bottomrule
\endfoot
1 & Revision churn (cand.\ 5), including the speaker-permutation ablation salvaged from cand.\ 2 & Seconds saved against retraction rate; band-flip rate at the 0.7 and 0.4 cuts & Cheapest by far. Needs only the capture path and one oracle rerun. No human annotation. Run it first to prove the harness. \\
\addlinespace
2 & Swapped-profile control (cand.\ 3 survivor) & Inter-role agreement, reported first; matched minus swapped effect size & The null test for the entire product thesis. If it fails, everything downstream must be reframed, so fail early. \\
\addlinespace
3 & Deadline-bounded ladder (cand.\ 1) & Fill correctness against p95 round trip; per-tier service fractions; silence precision & The lead claim, and the most engineering-heavy. Needs phase 0 items and a working peer protocol. \\
\addlinespace
4 & Confirmation-gated graph (cand.\ 4) & Edge precision against coverage; false-invalidation rate; write-path cost per confirmed edge & Needs the five-class hard-negative labelling, the most expensive annotation in the plan. \\
\addlinespace
5 & Abstention economics (cand.\ 6) & Three-region operating characteristic; per-role silent false-drop rate & Reuses phase 0 labels and phase 1 through 4 logs. Nearly free once the others have run. \\
\end{longtable}}

\subsection{Ablations reviewers will demand}

Four ablations are not optional, because each directly tests a claim the
system makes about itself. Omitting any one of them invites the reviewer to
assume the result would have been unfavourable.

\begin{itemize}
\item \textbf{Single-agent baseline.} Paper 3 finds that reported multi-agent
gains over single-agent baselines are often minimal. One well-prompted agent
over the same windows, producing the same six output types, is the control that
justifies six agents existing.
\item \textbf{Single-shared-summary baseline.} One summary handed to every
participant, scored against every per-role reference. This is the incumbent
product design and the thing the whole thesis claims to beat.
\item \textbf{Swapped and random profile controls.} Without these, any observed
per-participant difference is indistinguishable from prompt noise.
\item \textbf{Prompt-only contradiction detection.} The specific test of
whether graph traversal beats simply asking a model, which is the claim that
justifies the property graph.
\end{itemize}

\subsection{Reporting discipline}

Paper 4 sets the standard the agent literature now expects, and following it
costs little while pre-empting an entire class of objection.

Report a three axis frontier of quality, dollars and p95 latency, never quality
alone. State an admissibility rule for the live path and enforce it: retry,
temperature warming and serial escalation all buy accuracy with extra serial
calls, and all three are inadmissible inside a 1.5 second budget, so they must
not appear in a live-path baseline. Hold out the evaluation set from the golden
set that gates prompt changes, since fitting thresholds on the suite that also
gates deployment is the overfitting failure paper 4 names. Report
worst-participant quality alongside the mean, because theme 6 showed the mean
hides exactly the failure that matters. And report cost per participant-output
rather than per meeting, with cache hit rate broken out for shared stages
against per-role stages, which is the number that reveals whether the design
scales sublinearly in attendees or fragments the cache.

\section{Threats to validity and what reviewers will attack}

\subsection{The strongest objections}

\textbf{The corpus is staged, not real.} AMI meetings are scenario driven, with
assigned roles and a synthetic design task. If inter-role agreement turns out
high, that may be an artifact of staging rather than evidence that
personalization is unnecessary. Report the agreement statistic with this caveat
attached, and if any live Mytheus sessions exist, replicate on even a handful.

\textbf{The mechanism is engineering, not research.} This is the sharpest
attack on candidates 1 and 3, and in both cases it partly lands. Deadline
propagation and fallback chains are standard distributed systems practice, and
role-conditioned relevance appears to be patented. The response is the same in
both cases and it must be made explicitly rather than defensively: claim the
measurement, not the mechanism. What is unpublished is the operating
characteristic, not the idea.

\textbf{The evaluation is self-referential.} Several metrics are defined by
this project and measured on a substrate built by this project. That is
unavoidable when the gap is precisely that no substrate exists, but it must be
disclosed, and every new metric must be reported alongside a published anchor
metric on the same data so a reader can calibrate. Report R1-AUC beside the new
timeliness numbers, and published citation precision beside the utterance-level
decomposition.

\textbf{Annotation is the weak link.} Per-role relevance labels are subjective,
and the five-class decision-pair taxonomy needs genuine domain judgement.
Report chance-corrected agreement for every label set, publish the annotation
guidelines, and treat a low agreement figure as a finding rather than as a
problem to hide.

\textbf{The null may win.} If swapped profiles score close to matched, the
central thesis is weakened in print. This is the honest risk of running the
control at all, and it is also the reason the result is worth having, since
this is the only design in the corpus that can reject the null.

\subsection{How to frame the paper}

The framing that follows from all of the above is narrower than the system and
stronger than the pitch. Do not write a paper claiming a personalized
multi-agent meeting assistant, because the mechanism claims in that sentence
are each individually occupied. Write the paper that the empty quadrant
licenses.

\begin{takeaway}
\textbf{The defensible thesis.} Personalized, proactive meeting intelligence
has been built and patented, but never measured under the two constraints that
define it: a hard real-time budget on mutable speech, and a per-recipient
notion of correctness. This project supplies the missing substrate, the missing
controls, and the first operating characteristics. The contribution is that the
claims of this product category become falsifiable, and that on the evidence
here some of them do not survive.
\end{takeaway}

That framing has three properties worth noting. It survives the prior art that
refuted candidates 2 and 3, because it concedes the mechanism and claims the
measurement. It makes a negative result publishable, which protects the project
against its own strongest risk. And it puts the released artifacts, the
role-annotated slice, the hard negative taxonomy and the annotation guidelines,
at the centre, which is what makes a student project citable rather than merely
complete.
